\begin{helper}[Prefix-freeness of a ray]
Let $F$ be a front and let $s,t\in\mathsf{Ray}(F,n)$.  If $s$ is an
initial segment of $t$, then $s=t$.
\end{helper}
\begin{proof}
By \noderef{Ray}, every element of $s$ and $t$ is greater than $n$.
Suppose first that $s=t$.  Otherwise, by the definition of
\noderef{InitialSegment}, there is an $m\in t$ such that

\[
s=\{k\in t\mid k<m\}.
\]

Since $n<m$, inserting the common least element $n$ gives

\[
\{n\}\cup s
 =\{k\in\{n\}\cup t\mid k<m\}.
\]

Thus $\{n\}\cup s$ is an initial segment of $\{n\}\cup t$.  Both
inserted sets belong to $F$ by \noderef{Ray}, so prefix-freeness in
\noderef{Front} gives $\{n\}\cup s=\{n\}\cup t$.  Neither $s$ nor $t$
contains $n$, again by \noderef{Ray}; cancelling the common inserted
element yields $s=t$.
\end{proof}
